%%%%%%%%%%%%%%%%%%%%%%%%%%%%%%%%%%%%%%%%%%%%%%%%%%%%%%%%%%%%%
%% CATEGORY THEORY — Chapters 5–6
%% The Yoneda Lemma and Representability / Monads and Algebras
%%%%%%%%%%%%%%%%%%%%%%%%%%%%%%%%%%%%%%%%%%%%%%%%%%%%%%%%%%%%%

%%%%%%%%%%%%%%%%%%%%%%%%%%%%%%%%%%%%%%%%%%%%%%%%%%%%%%%%%%%%
%%    CHAPTER 5: THE YONEDA LEMMA AND REPRESENTABILITY    %%
%%%%%%%%%%%%%%%%%%%%%%%%%%%%%%%%%%%%%%%%%%%%%%%%%%%%%%%%%%%%

\chapter{The Yoneda Lemma and Representability}\label{ch:yoneda}
\minitoc

\vspace{1em}

The Yoneda lemma is, without exaggeration, the single most important
result in category theory.  It tells us that every object of a category
is completely determined by the totality of morphisms into (or out of) it.
More precisely, it establishes a bijection between natural
transformations out of a representable functor and the elements of the
target functor, and it shows that the Yoneda embedding is fully faithful.
This chapter develops the theory of presheaves, states and proves the
Yoneda lemma in its full generality, and explores representable functors
and the fundamental fact that every presheaf is a colimit of
representables.

\index{Yoneda lemma|(}
\index{representable functor|(}
\index{presheaf|(}

%% ============================================================
\section{Presheaves and the functor category}\label{sec:presheaves}
%% ============================================================

\begin{definition}[Presheaf]\label{def:presheaf}
\index{presheaf!definition}
\index{functor category!of presheaves}
Let $\mathcal{C}$ be a locally small category.
A \textbf{presheaf} on~$\mathcal{C}$ is a functor
\[
    F \colon \mathcal{C}^{\op} \longrightarrow \Set.
\]
The \textbf{category of presheaves} on~$\mathcal{C}$ is the functor category
\[
    \widehat{\mathcal{C}} \;=\; \Fun(\mathcal{C}^{\op},\,\Set)
    \;=\; \Set^{\mathcal{C}^{\op}}.
\]
\index{$\widehat{\mathcal{C}}$}
Morphisms in~$\widehat{\mathcal{C}}$ are natural transformations.
\end{definition}

\begin{remark}\label{rem:presheaf-covariant}
\index{copresheaf}
Dually, a \textbf{copresheaf} (or covariant presheaf) on~$\mathcal{C}$
is a functor $\mathcal{C}\to\Set$.  The category of copresheaves is
$\Fun(\mathcal{C},\Set)$.  Both viewpoints are important; the choice of
variance depends on the context.
\end{remark}

\begin{example}\label{ex:presheaf-topology}
\index{presheaf!on a topological space}
\index{open sets!poset of}
Let $X$ be a topological space and let $\mathcal{O}(X)$ be the poset of
open subsets ordered by inclusion, viewed as a category.
A presheaf on~$\mathcal{O}(X)$ in the sense above is a functor
$\mathcal{O}(X)^{\op}\to\Set$: it assigns to each open set~$U$ a
set~$F(U)$ and to each inclusion $V\subseteq U$ a \emph{restriction map}
$F(U)\to F(V)$, contravariantly.  This recovers the classical notion of
presheaf on a topological space.
\end{example}

\begin{example}\label{ex:presheaf-simplicial}
\index{simplicial set}
\index{$\mathbf{\Delta}$}
A \textbf{simplicial set} is a presheaf on the simplex category~$\mathbf{\Delta}$
(whose objects are the finite totally ordered sets $[n]=\{0,1,\ldots,n\}$
and whose morphisms are order-preserving maps).  Thus
$\mathsf{sSet} = \Fun(\mathbf{\Delta}^{\op},\Set)$.
\end{example}


%% ============================================================
\section{Representable functors}\label{sec:representable}
%% ============================================================

\begin{definition}[Covariant representable functor]\label{def:h-lower}
\index{representable functor!covariant}
\index{$h_A$}
\index{Hom functor!covariant}
For each object $A\in\mathcal{C}$, the \textbf{covariant representable
functor} (or \textbf{covariant Hom functor}) is
\[
    h_A \;=\; \Hom_{\mathcal{C}}(A,-) \;\colon\;
    \mathcal{C} \longrightarrow \Set.
\]
On objects: $h_A(B) = \Hom(A,B)$.
On morphisms: for $f\colon B\to B'$, the map
$h_A(f) = f_* = f\circ(-)\colon \Hom(A,B)\to\Hom(A,B')$
is post-composition with~$f$.
\end{definition}

\begin{definition}[Contravariant representable functor]\label{def:h-upper}
\index{representable functor!contravariant}
\index{$h^A$}
\index{Hom functor!contravariant}
For each object $A\in\mathcal{C}$, the \textbf{contravariant representable
functor} is
\[
    h^A \;=\; \Hom_{\mathcal{C}}(-,A) \;\colon\;
    \mathcal{C}^{\op} \longrightarrow \Set.
\]
On objects: $h^A(B) = \Hom(B,A)$.
On morphisms: for $f\colon B'\to B$ in~$\mathcal{C}$ (equivalently
$f\colon B\to B'$ in~$\mathcal{C}^{\op}$), the map
$h^A(f) = f^* = (-)\circ f\colon \Hom(B,A)\to\Hom(B',A)$
is pre-composition with~$f$.
\end{definition}

\begin{remark}\label{rem:representable-conventions}
\index{representable functor!convention}
Convention varies across the literature.  We follow the notation where
$h^A$ denotes the contravariant functor (a presheaf on~$\mathcal{C}$)
and $h_A$ denotes the covariant functor (a copresheaf).
Some authors write $\mathbf{y}(A)$ for $h^A$, where $\mathbf{y}$ is the
Yoneda embedding.
\end{remark}

\begin{definition}[Representable functor]\label{def:representable}
\index{representable functor!definition}
\index{representing object}
A functor $F\colon\mathcal{C}^{\op}\to\Set$ is \textbf{representable}
if there exists an object $A\in\mathcal{C}$ and a natural isomorphism
$F\cong h^A$.  The object~$A$ is called a \textbf{representing object}.
Dually, a functor $G\colon\mathcal{C}\to\Set$ is \textbf{representable}
if $G\cong h_A$ for some~$A$.
\end{definition}

\begin{notation}\label{not:hom-bifunctor}
\index{Hom functor!bifunctor}
\index{bifunctor!Hom}
The assignments $(A,B)\mapsto\Hom(A,B)$ assemble into a bifunctor
\[
    \Hom_{\mathcal{C}}(-,-) \;\colon\;
    \mathcal{C}^{\op}\times\mathcal{C} \longrightarrow \Set.
\]
Fixing the first variable yields $h_A = \Hom(A,-)$; fixing the second
yields $h^A = \Hom(-,A)$.
\end{notation}

\begin{example}\label{ex:representable-forgetful}
\index{forgetful functor!as representable}
\begin{enumerate}[label=(\roman*)]
\item The forgetful functor $U\colon\Grp\to\Set$ is representable:
    $U\cong h_{\mathbb{Z}} = \Hom_{\Grp}(\mathbb{Z},-)$, since a
    group homomorphism $\mathbb{Z}\to G$ is determined by the image
    of~$1$, hence by an element of the underlying set of~$G$.

\item The forgetful functor $U\colon\mathsf{Ring}\to\Set$ is
    representable: $U\cong\Hom_{\mathsf{Ring}}(\mathbb{Z}[x],-)$.

\item The forgetful functor $U\colon\Top\to\Set$ is representable:
    $U\cong\Hom_{\Top}(\{*\},-)$, where $\{*\}$ is the one-point space.

\item The functor $\mathsf{GL}_n\colon\mathsf{CRing}\to\Grp$ sending
    a commutative ring~$R$ to $\mathsf{GL}_n(R)$ is representable, with
    representing object $\mathbb{Z}[x_{ij},\det(x_{ij})^{-1}]$.
    \index{general linear group!as representable functor}
\end{enumerate}
\end{example}


%% ============================================================
\section{The Yoneda embedding}\label{sec:yoneda-embedding}
%% ============================================================

\begin{definition}[Yoneda embedding]\label{def:yoneda-embedding}
\index{Yoneda embedding}
\index{$\mathbf{y}$}
The \textbf{Yoneda embedding} is the functor
\[
    \mathbf{y} \;\colon\; \mathcal{C} \longrightarrow
    \Fun(\mathcal{C}^{\op},\Set) = \widehat{\mathcal{C}}
\]
defined on objects by $\mathbf{y}(A) = h^A = \Hom(-,A)$
and on morphisms by: for $f\colon A\to B$ in~$\mathcal{C}$,
$\mathbf{y}(f) = f_*\colon h^A\Rightarrow h^B$ is the natural transformation
whose component at $C\in\mathcal{C}$ is
\[
    (f_*)_C \;=\; f\circ(-) \;\colon\;
    \Hom(C,A) \longrightarrow \Hom(C,B).
\]
\end{definition}

\begin{lemma}\label{lem:yoneda-well-defined}
\index{Yoneda embedding!well-defined}
The map $\mathbf{y}(f) = f_*$ is indeed a natural transformation
$h^A\Rightarrow h^B$ for each morphism $f\colon A\to B$.
\end{lemma}

\begin{proof}
We must check naturality: for every $g\colon C'\to C$ in~$\mathcal{C}$,
the diagram
\[
\begin{tikzcd}
\Hom(C,A) \arrow[r, "f_*"] \arrow[d, "g^*"'] &
\Hom(C,B) \arrow[d, "g^*"] \\
\Hom(C',A) \arrow[r, "f_*"'] &
\Hom(C',B)
\end{tikzcd}
\]
commutes.  For $\varphi\in\Hom(C,A)$:
$g^*(f_*(\varphi)) = (f\circ\varphi)\circ g = f\circ(\varphi\circ g)
= f_*(g^*(\varphi))$.
\end{proof}

There is a dual version:

\begin{definition}[Co-Yoneda embedding]\label{def:co-yoneda-embedding}
\index{Yoneda embedding!covariant (co-Yoneda)}
The \textbf{co-Yoneda embedding} is the functor
\[
    \mathbf{y}_{\bullet} \;\colon\; \mathcal{C}^{\op}
    \longrightarrow \Fun(\mathcal{C},\Set)
\]
defined by $\mathbf{y}_{\bullet}(A) = h_A = \Hom(A,-)$.
\end{definition}


%% ============================================================
\section{The Yoneda lemma}\label{sec:yoneda-lemma}
%% ============================================================

We now state and prove the central result of this chapter.

\begin{theorem}[Yoneda Lemma]\label{thm:yoneda}
\index{Yoneda lemma!statement}
\index{Yoneda lemma!proof}
Let $\mathcal{C}$ be a locally small category, let $A\in\mathcal{C}$,
and let $F\colon\mathcal{C}^{\op}\to\Set$ be a presheaf.
There is a bijection
\[
    \Phi_{A,F} \;\colon\;
    \Nat(h^A,\,F) \;\xrightarrow{\;\sim\;}\; F(A)
\]
given by $\Phi_{A,F}(\alpha) = \alpha_A(\id_A)$, where
$\alpha\colon h^A\Rightarrow F$ is a natural transformation.
Moreover, this bijection is natural in both $A$ and~$F$.
\end{theorem}

\begin{proof}
The proof proceeds in four steps.

\medskip\noindent
\textbf{Step 1: Definition of the map~$\Phi$.}
Given a natural transformation $\alpha\colon h^A\Rightarrow F$,
we define
\[
    \Phi_{A,F}(\alpha) \;=\; \alpha_A(\id_A) \;\in\; F(A).
\]
This is well-defined since $\id_A\in h^A(A) = \Hom(A,A)$
and $\alpha_A\colon h^A(A)\to F(A)$.

\medskip\noindent
\textbf{Step 2: Definition of the inverse~$\Psi$.}
Given an element $x\in F(A)$, we define a natural transformation
$\Psi_{A,F}(x) = \alpha^x\colon h^A\Rightarrow F$ as follows.
For each $B\in\mathcal{C}$, the component
\[
    \alpha^x_B \;\colon\; h^A(B) = \Hom(B,A) \longrightarrow F(B)
\]
is defined by
\[
    \alpha^x_B(f) \;=\; F(f)(x) \quad\text{for } f\colon B\to A.
\]
Here $F(f)\colon F(A)\to F(B)$ is the action of the presheaf~$F$ on
the morphism $f$ (recall $F$ is contravariant on~$\mathcal{C}$, i.e.\
covariant on~$\mathcal{C}^{\op}$).

We must verify that $\alpha^x$ is a natural transformation.  Let
$g\colon C\to B$ be a morphism in~$\mathcal{C}$.  We need the following
diagram to commute:
\[
\begin{tikzcd}
h^A(B) \arrow[r, "\alpha^x_B"] \arrow[d, "h^A(g) = g^*"'] &
F(B) \arrow[d, "F(g)"] \\
h^A(C) \arrow[r, "\alpha^x_C"'] &
F(C)
\end{tikzcd}
\]
For $f\in h^A(B) = \Hom(B,A)$:
\begin{align*}
    F(g)(\alpha^x_B(f))
    &= F(g)(F(f)(x)) \\
    &= (F(g)\circ F(f))(x) \\
    &= F(f\circ g)(x) \quad\text{(since $F$ is a functor on $\mathcal{C}^{\op}$)}\\
    &= \alpha^x_C(f\circ g) \\
    &= \alpha^x_C(g^*(f)).
\end{align*}
Hence $\alpha^x$ is indeed natural.

\medskip\noindent
\textbf{Step 3: $\Phi$ and $\Psi$ are mutually inverse.}

\emph{$\Phi\circ\Psi = \id$:}
For $x\in F(A)$:
\[
    \Phi(\Psi(x)) = \Phi(\alpha^x) = \alpha^x_A(\id_A)
    = F(\id_A)(x) = \id_{F(A)}(x) = x.
\]

\emph{$\Psi\circ\Phi = \id$:}
For $\alpha\colon h^A\Rightarrow F$, let $x = \Phi(\alpha) = \alpha_A(\id_A)$.
We must show $\alpha^x = \alpha$, i.e.\ $\alpha^x_B = \alpha_B$ for
all $B\in\mathcal{C}$.  Let $f\in\Hom(B,A)$.  Consider the naturality
square of~$\alpha$ at $f\colon B\to A$:
\[
\begin{tikzcd}
h^A(A) \arrow[r, "\alpha_A"] \arrow[d, "f^*"'] &
F(A) \arrow[d, "F(f)"] \\
h^A(B) \arrow[r, "\alpha_B"'] &
F(B)
\end{tikzcd}
\]
Evaluating at $\id_A \in h^A(A)$:
\[
    \alpha_B(f^*(\id_A)) = F(f)(\alpha_A(\id_A)),
\]
i.e.\
\[
    \alpha_B(f) = F(f)(x) = \alpha^x_B(f).
\]
Since this holds for all $f\in\Hom(B,A)$, we have $\alpha_B = \alpha^x_B$
for all~$B$, hence $\alpha = \alpha^x = \Psi(\Phi(\alpha))$.

\medskip\noindent
\textbf{Step 4: Naturality in $A$ and~$F$.}

\emph{Naturality in~$A$:}
Let $g\colon A'\to A$ be a morphism.  We must show that the diagram
\[
\begin{tikzcd}
\Nat(h^A,F) \arrow[r, "\Phi_{A,F}"] \arrow[d, "g^*"'] &
F(A) \arrow[d, "F(g)"] \\
\Nat(h^{A'},F) \arrow[r, "\Phi_{A',F}"'] &
F(A')
\end{tikzcd}
\]
commutes, where $g^*(\alpha) = \alpha\circ\mathbf{y}(g) = \alpha\circ g_*$.
For $\alpha\in\Nat(h^A,F)$:
\begin{align*}
    \Phi_{A',F}(g^*(\alpha))
    &= \Phi_{A',F}(\alpha\circ g_*)
    = (\alpha\circ g_*)_{A'}(\id_{A'})
    = \alpha_{A'}(g\circ\id_{A'})
    = \alpha_{A'}(g).
\end{align*}
On the other hand, by the naturality of~$\alpha$ at~$g\colon A'\to A$:
\[
    F(g)(\alpha_A(\id_A)) = \alpha_{A'}(g) \quad\text{(from the naturality square)}.
\]
Hence $F(g)(\Phi_{A,F}(\alpha)) = \Phi_{A',F}(g^*(\alpha))$.

\emph{Naturality in~$F$:}
Let $\beta\colon F\Rightarrow G$ be a natural transformation of presheaves.
We must show that
\[
\begin{tikzcd}
\Nat(h^A,F) \arrow[r, "\Phi_{A,F}"] \arrow[d, "\beta_*"'] &
F(A) \arrow[d, "\beta_A"] \\
\Nat(h^A,G) \arrow[r, "\Phi_{A,G}"'] &
G(A)
\end{tikzcd}
\]
commutes, where $\beta_*(\alpha) = \beta\circ\alpha$.
For $\alpha\in\Nat(h^A,F)$:
\[
    \Phi_{A,G}(\beta_*(\alpha)) = \Phi_{A,G}(\beta\circ\alpha)
    = (\beta\circ\alpha)_A(\id_A) = \beta_A(\alpha_A(\id_A))
    = \beta_A(\Phi_{A,F}(\alpha)). \qedhere
\]
\end{proof}

\begin{remark}\label{rem:yoneda-key-idea}
\index{Yoneda lemma!key idea}
The essential insight of the Yoneda lemma is that a natural transformation
$\alpha\colon h^A\Rightarrow F$ is \emph{completely determined} by the
single element $\alpha_A(\id_A)\in F(A)$.  Naturality then forces
the value of~$\alpha$ at every other component.  This is sometimes
called the ``Yoneda trick'': evaluate at the identity.
\end{remark}

\begin{remark}[Dual Yoneda lemma]\label{rem:yoneda-dual}
\index{Yoneda lemma!dual}
The dual version states: for $A\in\mathcal{C}$ and a functor
$G\colon\mathcal{C}\to\Set$,
\[
    \Nat(h_A,\,G) \;\cong\; G(A),
\]
naturally in $A$ and~$G$.  The proof is formally dual.
\end{remark}


%% ============================================================
\section{Corollaries of the Yoneda lemma}\label{sec:yoneda-corollaries}
%% ============================================================

\begin{corollary}[Yoneda embedding is fully faithful]%
\label{cor:yoneda-ff}
\index{Yoneda embedding!fully faithful}
The Yoneda embedding $\mathbf{y}\colon\mathcal{C}\to\widehat{\mathcal{C}}$
is fully faithful.
\end{corollary}

\begin{proof}
For objects $A,B\in\mathcal{C}$, we must show that the map
\[
    \mathbf{y}_{A,B} \;\colon\;
    \Hom_{\mathcal{C}}(A,B) \longrightarrow
    \Nat(h^A,\,h^B) = \Hom_{\widehat{\mathcal{C}}}(\mathbf{y}(A),\mathbf{y}(B))
\]
is a bijection.  By the Yoneda lemma applied with $F = h^B$:
\[
    \Nat(h^A,\,h^B) \;\cong\; h^B(A) = \Hom(A,B).
\]
Moreover, the Yoneda bijection sends $\alpha\in\Nat(h^A,h^B)$ to
$\alpha_A(\id_A)\in\Hom(A,B)$.  For $\alpha = \mathbf{y}(f) = f_*$
(where $f\colon A\to B$), we get $(f_*)_A(\id_A) = f\circ\id_A = f$.
Thus the Yoneda bijection $\Nat(h^A,h^B)\xrightarrow{\sim}\Hom(A,B)$
is the inverse of~$\mathbf{y}_{A,B}$, so $\mathbf{y}_{A,B}$ is
a bijection.
\end{proof}

\begin{corollary}[Representables determine objects up to isomorphism]%
\label{cor:representable-iso}
\index{representable functor!determines object}
If $h^A \cong h^B$ as presheaves on~$\mathcal{C}$, then $A\cong B$
in~$\mathcal{C}$.  Dually, $h_A\cong h_B$ implies $A\cong B$.
\end{corollary}

\begin{proof}
Since $\mathbf{y}$ is fully faithful, it reflects isomorphisms
(see Exercise~\ref{exer:ff-reflects-iso}).
If $h^A\cong h^B$ in $\widehat{\mathcal{C}}$, let
$\alpha\colon h^A\xrightarrow{\sim} h^B$ be a natural isomorphism.
Since $\mathbf{y}$ is full, there exists $f\colon A\to B$ with
$\mathbf{y}(f)=\alpha$.  Similarly, the inverse $\alpha^{-1}$ equals
$\mathbf{y}(g)$ for some $g\colon B\to A$.  Since $\mathbf{y}$ is faithful:
\[
    \mathbf{y}(g\circ f) = \mathbf{y}(g)\circ\mathbf{y}(f)
    = \alpha^{-1}\circ\alpha = \id_{h^A} = \mathbf{y}(\id_A)
    \implies g\circ f = \id_A.
\]
Similarly $f\circ g = \id_B$, so $f$ is an isomorphism.
\end{proof}

\begin{corollary}[Uniqueness of representing objects]%
\label{cor:rep-unique}
\index{representing object!uniqueness}
If a functor $F\colon\mathcal{C}^{\op}\to\Set$ is represented by both
$A$ and~$A'$, then $A\cong A'$ in~$\mathcal{C}$.
\end{corollary}

\begin{proof}
If $F\cong h^A$ and $F\cong h^{A'}$, then $h^A\cong h^{A'}$, hence
$A\cong A'$ by Corollary~\ref{cor:representable-iso}.
\end{proof}

\begin{corollary}\label{cor:yoneda-universal-element}
\index{universal element}
Let $F\colon\mathcal{C}^{\op}\to\Set$.  Then $F$ is representable
(i.e.\ $F\cong h^A$ for some~$A$) if and only if there exists an object
$A\in\mathcal{C}$ and an element $u\in F(A)$---called the
\textbf{universal element}---such that for every $B\in\mathcal{C}$, the map
\[
    \Hom(B,A) \longrightarrow F(B), \quad f\longmapsto F(f)(u)
\]
is a bijection.
\end{corollary}

\begin{proof}
By the Yoneda lemma, a natural isomorphism $\alpha\colon h^A\xrightarrow{\sim} F$
corresponds to the element $u = \alpha_A(\id_A)\in F(A)$.
The component $\alpha_B\colon\Hom(B,A)\to F(B)$ sends $f$ to
$F(f)(u)$ (this is the formula from Step~2 of the proof of the
Yoneda lemma).  Conversely, any $u\in F(A)$ such that $f\mapsto F(f)(u)$
is bijective for all~$B$ defines a natural isomorphism
$h^A\xrightarrow{\sim}F$.
\end{proof}


%% ============================================================
\section{Every presheaf is a colimit of representables}%
\label{sec:colim-representables}
%% ============================================================

\begin{definition}[Category of elements]\label{def:category-of-elements}
\index{category of elements}
\index{$\int F$}
\index{Grothendieck construction!for presheaves}
Let $F\colon\mathcal{C}^{\op}\to\Set$ be a presheaf.
The \textbf{category of elements} of~$F$, denoted $\int F$
(or $\mathcal{C}/F$ or $\mathrm{el}(F)$), has:
\begin{itemize}
\item \textbf{Objects:} pairs $(C,x)$ with $C\in\mathcal{C}$ and $x\in F(C)$.
\item \textbf{Morphisms:} a morphism $(C,x)\to(C',x')$ is a morphism
    $f\colon C\to C'$ in~$\mathcal{C}$ such that $F(f)(x') = x$
    (note the contravariance).
\end{itemize}
There is a canonical projection functor
$\pi_F\colon\int F\to\mathcal{C}$, $(C,x)\mapsto C$.
\end{definition}

\begin{example}\label{ex:elements-representable}
\index{category of elements!of representable}
If $F = h^A$, then an object of $\int h^A$ is a pair $(C,f)$ with
$f\in\Hom(C,A)$, i.e.\ a morphism $f\colon C\to A$.
A morphism $(C,f)\to(C',f')$ is $g\colon C\to C'$ with $f = f'\circ g$.
Thus $\int h^A \cong \mathcal{C}/A$, the slice category over~$A$.
\end{example}

\begin{theorem}[Every presheaf is a colimit of representables]%
\label{thm:presheaf-colim-rep}
\index{presheaf!as colimit of representables}
\index{colimit!of representables}
\index{representable functor!presheaf as colimit of}
Let $F\colon\mathcal{C}^{\op}\to\Set$ be a presheaf.  Then $F$ is the
colimit, in the presheaf category~$\widehat{\mathcal{C}}$, of the
diagram of representable presheaves indexed by the category of elements:
\[
    F \;\cong\; \varinjlim_{(C,x)\,\in\,\int F}\; h^C.
\]
More precisely, the composite functor
$\int F \xrightarrow{\pi_F} \mathcal{C} \xrightarrow{\mathbf{y}}
\widehat{\mathcal{C}}$
has colimit~$F$.
\end{theorem}

\begin{proof}
We construct a cocone and verify the universal property.
For each $(C,x)\in\int F$, the element $x\in F(C)$ corresponds, by the
Yoneda lemma, to a natural transformation
$\alpha^x\colon h^C\Rightarrow F$
defined by $\alpha^x_B(f) = F(f)(x)$ for $f\in\Hom(B,C)$.
This gives a cocone $\{\alpha^x\colon h^C\to F\}_{(C,x)\in\int F}$.

We verify compatibility: if $g\colon(C,x)\to(C',x')$ is a morphism in
$\int F$ (so $g\colon C\to C'$ and $F(g)(x') = x$), then for all
$B\in\mathcal{C}$ and $f\in\Hom(B,C)$:
\[
    \alpha^{x'}_B(g\circ f) = F(g\circ f)(x') = F(f)(F(g)(x'))
    = F(f)(x) = \alpha^x_B(f),
\]
so $\alpha^{x'}\circ g_* = \alpha^x$ as required.

Now let $\{\beta^{(C,x)}\colon h^C\to G\}_{(C,x)\in\int F}$ be any
cocone over the diagram.  We define a natural transformation
$\gamma\colon F\Rightarrow G$ by
\[
    \gamma_B(y) \;=\; \beta^{(B,y)}_B(\id_B) \quad\text{for } y\in F(B).
\]
One checks naturality of~$\gamma$ using the cocone compatibility, and
uniqueness follows because any $\gamma$ satisfying
$\gamma\circ\alpha^x = \beta^{(C,x)}$ for all $(C,x)$ must satisfy
$\gamma_C(x) = \gamma_C(\alpha^x_C(\id_C)) = \beta^{(C,x)}_C(\id_C)$.
\end{proof}

\begin{remark}\label{rem:density}
\index{density theorem}
\index{Yoneda embedding!dense}
Theorem~\ref{thm:presheaf-colim-rep} is sometimes called the
\textbf{density theorem} or the statement that the Yoneda embedding
is \textbf{dense}.  It is fundamental in the theory of Kan extensions
and in the study of locally presentable categories.
\end{remark}


%% ============================================================
\section{Applications of the Yoneda lemma}\label{sec:yoneda-applications}
%% ============================================================

\begin{proposition}[Natural transformations via representables]%
\label{prop:nat-trans-representable}
\index{natural transformation!between representable functors}
For any functor $F\colon\mathcal{C}\to\mathcal{D}$ between locally small
categories, and objects $A,B\in\mathcal{C}$, there is a bijection
\[
    \Nat(h^{F(A)},\, h^{F(B)}) \;\cong\; \Hom_{\mathcal{D}}(F(A),F(B)).
\]
In particular, if $F$ is fully faithful, then it induces a fully faithful
functor on presheaf categories.
\end{proposition}

\begin{proof}
This is the Yoneda lemma applied to~$\mathcal{D}$ with
$G = h^{F(B)}$:
$\Nat(h^{F(A)},h^{F(B)})\cong h^{F(B)}(F(A)) = \Hom(F(A),F(B))$.
\end{proof}

\begin{proposition}[Characterisation of limits via representables]%
\label{prop:limit-representable}
\index{limit!via representable functor}
Let $D\colon\mathcal{J}\to\mathcal{C}$ be a diagram.  An object
$L\in\mathcal{C}$ together with a cone $(\pi_j\colon L\to D_j)$ is a
limit of~$D$ if and only if for every $C\in\mathcal{C}$, the canonical map
\[
    \Hom(C,L) \longrightarrow \varprojlim_{j\in\mathcal{J}}\,\Hom(C,D_j)
\]
is a bijection.  Equivalently, $h^L\cong\varprojlim_j\, h^{D_j}$
as presheaves, i.e.\ limits in~$\mathcal{C}$ are computed pointwise
by the Yoneda embedding.
\end{proposition}

\begin{proof}
By definition, $L = \varprojlim D$ means
$\Hom(C,L)\cong\mathrm{Cone}(C,D)$ naturally in~$C$.
The set of cones from~$C$ to~$D$ is precisely
$\varprojlim_j\Hom(C,D_j)$ (where the limit is in~$\Set$).
Naturality in~$C$ means $h^L\cong\varprojlim_j h^{D_j}$ as presheaves.
Since the Yoneda embedding is fully faithful, this limit in the
presheaf category reflects to a limit in~$\mathcal{C}$.
\end{proof}


%% ============================================================
\section{Exercises}\label{sec:yoneda-exercises}
%% ============================================================

\begin{exercise}\label{exer:ff-reflects-iso}
\index{fully faithful!reflects isomorphisms}
Show that a fully faithful functor reflects isomorphisms: if
$F\colon\mathcal{C}\to\mathcal{D}$ is fully faithful and $F(f)$
is an isomorphism, then $f$ is an isomorphism.
\end{exercise}

\begin{exercise}\label{exer:yoneda-explicit}
\index{Yoneda lemma!explicit example}
Let $\mathcal{C} = \mathsf{Vect}_k$ (vector spaces over a field~$k$).
\begin{enumerate}[label=(\alph*)]
\item Show that $\Nat(h^V,h^W)\cong\Hom(V,W)$ for all vector
    spaces $V,W$ by tracing through the Yoneda bijection.
\item Let $F(V) = V\otimes V$ (a covariant functor).  Compute
    $\Nat(h_k,\,F)$ and identify it with~$F(k) = k$.
\end{enumerate}
\end{exercise}

\begin{exercise}\label{exer:group-as-presheaf}
\index{group!as presheaf}
Let $\mathcal{C}$ be a category with finite products.  A \textbf{group
object} in~$\mathcal{C}$ is an object~$G$ with morphisms
$m\colon G\times G\to G$, $e\colon 1\to G$, $\iota\colon G\to G$
satisfying the usual diagrams.  Show that~$G$ is a group object if and
only if $h^G = \Hom(-,G)\colon\mathcal{C}^{\op}\to\Set$ factors through
$\Grp\to\Set$.
\emph{Hint:} Use the Yoneda lemma to translate element-free diagrams
to element-wise conditions.
\end{exercise}

\begin{exercise}\label{exer:yoneda-monoid}
\index{Yoneda lemma!for monoids}
Let $M$ be a monoid viewed as a one-object category~$\mathcal{B}_M$.
\begin{enumerate}[label=(\alph*)]
\item Show that a presheaf on~$\mathcal{B}_M$ is the same as a
    right $M$-set.
\item Identify the representable presheaf $h^{*}$ (where $*$ is
    the unique object) with $M$ acting on itself by right multiplication.
\item Formulate and verify the Yoneda lemma in this case: for every
    right $M$-set~$S$, the set of $M$-equivariant maps $M\to S$
    is in bijection with~$S$.
\end{enumerate}
\end{exercise}

\begin{exercise}\label{exer:subfunctor}
\index{subfunctor}
\index{sieve}
Let $F\colon\mathcal{C}^{\op}\to\Set$ be a presheaf and let
$S\subseteq F$ be a subfunctor ($S(C)\subseteq F(C)$ for all~$C$,
compatible with restriction maps).  Show that every subfunctor
of~$h^A$ corresponds to a \textbf{sieve} on~$A$, i.e.\ a set~$S$ of
morphisms with codomain~$A$ such that if $f\in S$ and $g$ is composable,
then $f\circ g\in S$.
\end{exercise}

\begin{exercise}\label{exer:cayley-yoneda}
\index{Cayley's theorem!via Yoneda}
\index{Yoneda lemma!Cayley's theorem as special case}
Show that Cayley's theorem in group theory (every group embeds in a
symmetric group) is a special case of the Yoneda lemma.
\emph{Hint:} Apply the Yoneda embedding to the one-object category
associated to a group.
\end{exercise}

\begin{exercise}\label{exer:presheaf-colim-rep}
\index{presheaf!as colimit of representables}
Fill in the remaining details of the proof of
Theorem~\ref{thm:presheaf-colim-rep}: verify that $\gamma$ as defined
there is indeed natural.
\end{exercise}

\begin{exercise}\label{exer:yoneda-enriched}
\index{Yoneda lemma!enriched}
\index{enriched category}
Let $\mathcal{V}$ be a complete, cocomplete, symmetric monoidal closed
category.  Formulate (without proof) the enriched Yoneda lemma for a
$\mathcal{V}$-enriched category~$\mathcal{C}$: for a
$\mathcal{V}$-functor $F\colon\mathcal{C}^{\op}\to\mathcal{V}$ and
$A\in\mathcal{C}$,
\[
    [\mathcal{C}^{\op},\mathcal{V}](h^A,F) \;\cong\; F(A)
\]
in~$\mathcal{V}$, where the left-hand side is the $\mathcal{V}$-object
of $\mathcal{V}$-natural transformations.
\end{exercise}

\index{Yoneda lemma|)}
\index{representable functor|)}
\index{presheaf|)}


%%%%%%%%%%%%%%%%%%%%%%%%%%%%%%%%%%%%%%%%%%%%%%%%%%%%%%%%%%%%
%%        CHAPTER 6: MONADS AND ALGEBRAS                  %%
%%%%%%%%%%%%%%%%%%%%%%%%%%%%%%%%%%%%%%%%%%%%%%%%%%%%%%%%%%%%

\chapter{Monads and Algebras}\label{ch:monads}
\minitoc

\vspace{1em}

Monads are one of the most versatile concepts in category theory,
appearing in algebra, topology, logic, and computer science.
A monad on a category captures the essence of algebraic structure:
it packages together a ``free construction'' and the ``equations''
that the resulting algebraic objects must satisfy.
In this chapter we define monads, show how they arise from adjunctions,
study the categories of algebras over a monad (the Eilenberg--Moore and
Kleisli categories), and prove Beck's monadicity theorem, which
characterises those adjunctions that arise from monads.

\index{monad|(}

%% ============================================================
\section{Monads}\label{sec:monad-def}
%% ============================================================

\begin{definition}[Monad]\label{def:monad}
\index{monad!definition}
\index{unit!of a monad}
\index{multiplication!of a monad}
A \textbf{monad} on a category~$\mathcal{C}$ is a triple
$(T,\eta,\mu)$ where:
\begin{itemize}
\item $T\colon\mathcal{C}\to\mathcal{C}$ is an endofunctor,
\item $\eta\colon\Id_{\mathcal{C}}\Rightarrow T$ is a natural
    transformation called the \textbf{unit},
\item $\mu\colon T^2\Rightarrow T$ is a natural transformation called
    the \textbf{multiplication},
\end{itemize}
subject to the following diagrams commuting for every object
$A\in\mathcal{C}$:

\textbf{Associativity:}
\[
\begin{tikzcd}
T^3A \arrow[r, "T\mu_A"] \arrow[d, "\mu_{TA}"'] &
T^2A \arrow[d, "\mu_A"] \\
T^2A \arrow[r, "\mu_A"'] &
TA
\end{tikzcd}
\]

\textbf{Unit laws:}
\[
\begin{tikzcd}
TA \arrow[r, "T\eta_A"] \arrow[dr, equal] &
T^2A \arrow[d, "\mu_A"] &
TA \arrow[l, "\eta_{TA}"'] \arrow[dl, equal] \\
& TA &
\end{tikzcd}
\]
Equivalently, $\mu\circ T\mu = \mu\circ\mu T$ and
$\mu\circ T\eta = \mu\circ\eta T = \id_T$.
\end{definition}

\begin{remark}\label{rem:monad-as-monoid}
\index{monad!as monoid in endofunctor category}
\index{monoid!in monoidal category}
A monad is a \textbf{monoid} in the monoidal category
$(\mathrm{End}(\mathcal{C}),\circ,\Id)$ of endofunctors with
composition as the monoidal product.  The unit~$\eta$ is the monoid
unit and the multiplication~$\mu$ is the monoid multiplication.
This perspective is extremely useful.
\end{remark}


%% ============================================================
\section{Monads from adjunctions}\label{sec:monad-adj}
%% ============================================================

\begin{proposition}[Every adjunction gives a monad]%
\label{prop:adj-monad}
\index{monad!from adjunction}
\index{adjunction!associated monad}
Let $F\dashv G$ be an adjunction with
$F\colon\mathcal{C}\to\mathcal{D}$, $G\colon\mathcal{D}\to\mathcal{C}$,
unit $\eta\colon\Id_{\mathcal{C}}\Rightarrow GF$, and counit
$\varepsilon\colon FG\Rightarrow\Id_{\mathcal{D}}$.
Then $(T,\eta,\mu)$ is a monad on~$\mathcal{C}$, where $T = GF$ and
$\mu = G\varepsilon F\colon GFGF\Rightarrow GF$.
\end{proposition}

\begin{proof}
\textit{Associativity:}
We must show $\mu\circ T\mu = \mu\circ\mu T$, i.e.\
$G\varepsilon F\circ GF\cdot G\varepsilon F
= G\varepsilon F\circ G\varepsilon F\cdot GF$.
Equivalently, for each~$A$:
\begin{align*}
    \mu_A \circ T\mu_A
    &= G\varepsilon_{FA}\circ GFG\varepsilon_{FA}
    = G(\varepsilon_{FA}\circ FG\varepsilon_{FA}), \\
    \mu_A \circ \mu_{TA}
    &= G\varepsilon_{FA}\circ G\varepsilon_{FGF A}
    = G(\varepsilon_{FA}\circ\varepsilon_{FGFA}).
\end{align*}
By naturality of~$\varepsilon$ at $\varepsilon_{FA}\colon FGFA\to FA$:
$\varepsilon_{FA}\circ FG\varepsilon_{FA} = \varepsilon_{FA}\circ\varepsilon_{FGFA}$.
Hence the two expressions agree.

\textit{Left unit law:} We show $\mu\circ\eta T = \id_T$.
For each~$A$:
$\mu_A\circ\eta_{TA} = G\varepsilon_{FA}\circ\eta_{GFA}
= \id_{GFA} = \id_{TA}$,
by one of the triangle identities ($G\varepsilon\circ\eta G = \id_G$
applied to $FA$).

\textit{Right unit law:} We show $\mu\circ T\eta = \id_T$.
For each~$A$:
$\mu_A\circ T\eta_A = G\varepsilon_{FA}\circ GF\eta_A
= G(\varepsilon_{FA}\circ F\eta_A) = G(\id_{FA}) = \id_{TA}$,
by the other triangle identity ($\varepsilon F\circ F\eta = \id_F$
applied to~$A$).
\end{proof}

\begin{remark}\label{rem:comonad}
\index{comonad}
\index{comonad!from adjunction}
Dually, the adjunction $F\dashv G$ also gives a \textbf{comonad}
$(S,\varepsilon,\delta)$ on~$\mathcal{D}$, where $S = FG$,
the counit is $\varepsilon\colon FG\Rightarrow\Id_{\mathcal{D}}$,
and the comultiplication is $\delta = F\eta G\colon FG\Rightarrow FGFG$.
\end{remark}


%% ============================================================
\section{Examples of monads}\label{sec:monad-examples}
%% ============================================================

\begin{example}[Free monoid monad]\label{ex:monad-free-monoid}
\index{monad!free monoid}
\index{free monoid}
The free-forgetful adjunction $F\dashv U\colon\mathsf{Mon}\to\Set$
gives rise to the monad $T = UF$ on~$\Set$.  Explicitly:
\[
    T(X) \;=\; \coprod_{n\geq 0} X^n \;=\;
    \{(x_1,\ldots,x_n) : n\geq 0,\, x_i\in X\},
\]
the set of finite words on the alphabet~$X$.
The unit $\eta_X\colon X\to T(X)$ sends $x$ to the one-letter
word~$(x)$.  The multiplication $\mu_X\colon T^2(X)\to T(X)$ is
``flattening'' (concatenation): it sends a word of words to a single
word by removing parentheses.  The monad axioms correspond to the
associativity of concatenation and the fact that the empty word is
a unit.
\end{example}

\begin{example}[Free group monad]\label{ex:monad-free-group}
\index{monad!free group}
\index{free group}
The adjunction $F\dashv U\colon\Grp\to\Set$ gives a monad~$T$
where $T(X)$ is the underlying set of the free group on~$X$, i.e.\
the set of reduced words in $X\cup X^{-1}$.
\end{example}

\begin{example}[Free $R$-module monad]\label{ex:monad-free-module}
\index{monad!free module}
\index{free module}
For a ring~$R$, the adjunction $F\dashv U\colon R\text{-}\mathsf{Mod}\to\Set$
gives a monad with
\[
    T(X) \;=\; \bigoplus_{x\in X} R \;=\;
    \Bigl\{\sum_{x\in X} r_x\cdot x :
    r_x\in R,\;\text{finitely many nonzero}\Bigr\}.
\]
\end{example}

\begin{example}[Power set monad]\label{ex:monad-powerset}
\index{monad!power set}
\index{power set monad}
The \textbf{power set monad} $\mathcal{P}$ on $\Set$ is defined by
$\mathcal{P}(X) = 2^X$ (the power set), with
$\eta_X(x) = \{x\}$ and $\mu_X(\mathcal{A}) = \bigcup_{A\in\mathcal{A}} A$
for $\mathcal{A}\in\mathcal{P}(\mathcal{P}(X))$.
This monad arises from the adjunction between $\Set$ and the category
of complete sup-lattices.
\end{example}

\begin{example}[Maybe monad]\label{ex:monad-maybe}
\index{monad!maybe}
\index{maybe monad}
The functor $T(X) = X\sqcup\{*\}$ on $\Set$, with $\eta_X$ the
inclusion $X\hookrightarrow X\sqcup\{*\}$ and $\mu_X$ collapsing
the two copies of~$*$, is a monad.  In computer science, this is the
\textbf{Maybe} (or \textbf{option}) monad, modelling computations that
may fail.
\end{example}

\begin{example}[Continuation monad]\label{ex:monad-continuation}
\index{monad!continuation}
\index{continuation monad}
Fix a set~$R$.  The functor $T(X) = R^{R^X}$ on $\Set$ carries a
monad structure called the \textbf{continuation monad}.
In functional programming, it arises from the self-adjunction
$(-)^R\dashv(-)^R\colon\Set^{\op}\to\Set$.
\end{example}

\begin{example}[Monad from a closure operator]\label{ex:monad-closure}
\index{monad!closure operator}
\index{closure operator}
Let $(P,\leq)$ be a poset viewed as a category.  A monad on~$P$ is a
\textbf{closure operator}: a monotone map $c\colon P\to P$ with
$x\leq c(x)$ (unit) and $c(c(x))\leq c(x)$ (multiplication), which
combined with monotonicity gives $c(c(x)) = c(x)$ (idempotence).
\end{example}


%% ============================================================
\section{Eilenberg--Moore algebras}\label{sec:em-algebras}
%% ============================================================

\begin{definition}[Algebra over a monad]\label{def:em-algebra}
\index{Eilenberg--Moore algebra}
\index{monad!algebra}
\index{$T$-algebra}
Let $(T,\eta,\mu)$ be a monad on~$\mathcal{C}$.
A \textbf{$T$-algebra} (or \textbf{Eilenberg--Moore algebra}) is a pair
$(A,a)$ where $A\in\mathcal{C}$ is an object and
$a\colon TA\to A$ is a morphism (the \textbf{structure map}) such that:

\textbf{Unit:}
$a\circ\eta_A = \id_A$.
\[
\begin{tikzcd}
A \arrow[r, "\eta_A"] \arrow[dr, equal] &
TA \arrow[d, "a"] \\
& A
\end{tikzcd}
\]

\textbf{Associativity:}
$a\circ\mu_A = a\circ Ta$.
\[
\begin{tikzcd}
T^2A \arrow[r, "\mu_A"] \arrow[d, "Ta"'] &
TA \arrow[d, "a"] \\
TA \arrow[r, "a"'] &
A
\end{tikzcd}
\]
\end{definition}

\begin{definition}[Morphism of algebras]\label{def:em-morphism}
\index{Eilenberg--Moore algebra!morphism}
\index{$T$-algebra!morphism}
A \textbf{morphism of $T$-algebras} $f\colon(A,a)\to(B,b)$ is a
morphism $f\colon A\to B$ in~$\mathcal{C}$ such that $f\circ a = b\circ Tf$:
\[
\begin{tikzcd}
TA \arrow[r, "Tf"] \arrow[d, "a"'] &
TB \arrow[d, "b"] \\
A \arrow[r, "f"'] &
B
\end{tikzcd}
\]
\end{definition}

\begin{definition}[Eilenberg--Moore category]\label{def:em-category}
\index{Eilenberg--Moore category}
\index{$\mathcal{C}^T$}
The \textbf{Eilenberg--Moore category} $\mathcal{C}^T$ has:
\begin{itemize}
\item \textbf{Objects:} $T$-algebras $(A,a)$,
\item \textbf{Morphisms:} morphisms of $T$-algebras,
\item \textbf{Composition and identities:} inherited from~$\mathcal{C}$.
\end{itemize}
There is a forgetful functor $U^T\colon\mathcal{C}^T\to\mathcal{C}$,
$(A,a)\mapsto A$, which has a left adjoint
$F^T\colon\mathcal{C}\to\mathcal{C}^T$ defined by
$F^T(A) = (TA,\mu_A)$.
\end{definition}

\begin{proposition}\label{prop:em-adjunction}
\index{Eilenberg--Moore category!adjunction}
The pair $F^T\dashv U^T$ is an adjunction, and the monad arising
from this adjunction (via Proposition~\ref{prop:adj-monad}) is
precisely~$(T,\eta,\mu)$.
\end{proposition}

\begin{proof}
\textit{$F^T(A) = (TA,\mu_A)$ is a $T$-algebra:}
The unit axiom is $\mu_A\circ\eta_{TA} = \id_{TA}$ (left unit of the
monad).  The associativity axiom is $\mu_A\circ\mu_{TA} = \mu_A\circ T\mu_A$
(monad associativity).

\textit{Adjunction:} For $(A,a)\in\mathcal{C}^T$ and $B\in\mathcal{C}$,
define
\[
    \Phi\colon\Hom_{\mathcal{C}^T}(F^TB,\,(A,a))
    \longrightarrow \Hom_{\mathcal{C}}(B,\,U^T(A,a) = A)
\]
by $\Phi(f) = f\circ\eta_B$.  The inverse is
$\Psi(g) = a\circ Tg$ for $g\colon B\to A$.
One verifies that $\Psi(g)$ is a $T$-algebra morphism and that
$\Phi,\Psi$ are mutually inverse using the algebra axioms and monad axioms.

\textit{Recovery of the monad:}
$U^TF^T(A) = U^T(TA,\mu_A) = TA = T(A)$, so $U^TF^T = T$.
The unit of the adjunction $F^T\dashv U^T$ is $\eta$.
The counit at $(A,a)$ is the algebra map
$\varepsilon_{(A,a)} = a\colon(T A,\mu_A)\to(A,a)$,
which is a $T$-algebra morphism by the associativity axiom for~$(A,a)$.
The induced multiplication is $U^T\varepsilon_{F^T}\cdot\eta = \mu$,
as desired.
\end{proof}

\begin{example}\label{ex:em-algebras}
\index{Eilenberg--Moore algebra!examples}
\begin{enumerate}[label=(\roman*)]
\item \textbf{Free monoid monad} (Example~\ref{ex:monad-free-monoid}):
    $T$-algebras are monoids.  A structure map $a\colon T(X)\to X$
    takes a word $(x_1,\ldots,x_n)$ and produces a single element:
    this is the multiplication.  The algebra axioms force associativity
    and unitality.
    \index{monoid!as Eilenberg--Moore algebra}

\item \textbf{Free group monad} (Example~\ref{ex:monad-free-group}):
    $T$-algebras are groups.  More precisely,
    $\Set^T\simeq\Grp$.
    \index{group!as Eilenberg--Moore algebra}

\item \textbf{Free module monad} (Example~\ref{ex:monad-free-module}):
    $T$-algebras are $R$-modules, i.e.\ $\Set^T\simeq R\text{-}\mathsf{Mod}$.
    \index{module!as Eilenberg--Moore algebra}

\item \textbf{Power set monad} (Example~\ref{ex:monad-powerset}):
    $T$-algebras are complete sup-lattices with sup-preserving maps
    as morphisms.
    \index{complete lattice!as Eilenberg--Moore algebra}

\item \textbf{Closure operators} (Example~\ref{ex:monad-closure}):
    the algebras for a closure operator $c$ on a poset~$P$ are exactly
    the \textbf{closed elements}, i.e.\ those $x$ with $c(x) = x$.
    \index{closed element}
\end{enumerate}
\end{example}


%% ============================================================
\section{Free algebras}\label{sec:free-algebras}
%% ============================================================

\begin{definition}[Free $T$-algebra]\label{def:free-algebra}
\index{free algebra}
\index{$T$-algebra!free}
A $T$-algebra $(A,a)$ is \textbf{free} if it is isomorphic to
$F^T(X) = (TX,\mu_X)$ for some $X\in\mathcal{C}$.
\end{definition}

\begin{proposition}\label{prop:free-algebra-universal}
\index{free algebra!universal property}
The free $T$-algebra on~$X$ satisfies the universal property:
for every $T$-algebra $(A,a)$ and every morphism $f\colon X\to A$
in~$\mathcal{C}$, there is a unique $T$-algebra morphism
$\bar{f}\colon(TX,\mu_X)\to(A,a)$ with $\bar{f}\circ\eta_X = f$.
Explicitly, $\bar{f} = a\circ Tf$.
\end{proposition}

\begin{proof}
This is simply the adjunction $F^T\dashv U^T$ applied to the pair
$(X,(A,a))$: $\Hom_{\mathcal{C}^T}(F^TX,(A,a))\cong\Hom_{\mathcal{C}}(X,A)$.
\end{proof}


%% ============================================================
\section{The Kleisli category}\label{sec:kleisli}
%% ============================================================

\begin{definition}[Kleisli category]\label{def:kleisli}
\index{Kleisli category}
\index{$\mathcal{C}_T$}
Let $(T,\eta,\mu)$ be a monad on~$\mathcal{C}$.
The \textbf{Kleisli category} $\mathcal{C}_T$ has:
\begin{itemize}
\item \textbf{Objects:} the same as~$\mathcal{C}$,
\item \textbf{Morphisms:} $\Hom_{\mathcal{C}_T}(A,B) = \Hom_{\mathcal{C}}(A,TB)$,
\item \textbf{Identity:} $\id_A = \eta_A\colon A\to TA$,
\item \textbf{Composition:} given $f\colon A\to TB$ and $g\colon B\to TC$
    (in~$\mathcal{C}$), their Kleisli composite is
    \[
        g\circ_T f \;=\; \mu_C\circ Tg\circ f \;\colon\; A\to TC.
    \]
\end{itemize}
\end{definition}

\begin{proposition}\label{prop:kleisli-category}
\index{Kleisli category!well-defined}
The Kleisli category~$\mathcal{C}_T$ is well-defined, i.e.\
composition is associative and $\eta_A$ is a two-sided identity.
\end{proposition}

\begin{proof}
\textit{Left identity:}
$\id_B\circ_T f = \mu_B\circ T\eta_B\circ f = \id_{TB}\circ f = f$
(right unit law of the monad).

\textit{Right identity:}
$f\circ_T\id_A = \mu_B\circ Tf\circ\eta_A = f$
(naturality of~$\eta$: $Tf\circ\eta_A = \eta_{TB}\circ f$, then
$\mu_B\circ\eta_{TB} = \id_{TB}$).

\textit{Associativity:}
Let $f\colon A\to TB$, $g\colon B\to TC$, $h\colon C\to TD$.  Then:
\begin{align*}
    h\circ_T(g\circ_T f)
    &= \mu_D\circ Th\circ\mu_C\circ Tg\circ f, \\
    (h\circ_T g)\circ_T f
    &= \mu_D\circ T(\mu_D\circ Th\circ g)\circ f
    = \mu_D\circ T\mu_D\circ T^2h\circ Tg\circ f.
\end{align*}
By naturality of~$\mu$ at~$h$: $\mu_D\circ Th = \mu_D\circ Th$.
More precisely, we use $\mu_D\circ T\mu_D = \mu_D\circ\mu_{TD}$
(associativity of the monad) and naturality of~$\mu$ at~$h$:
$\mu_D\circ Th = Th \circ \mu_C$ is not quite right; let us be
more careful.  By naturality of~$\mu$ at $h\colon C\to TD$,
$\mu_D\circ T^2h = Th\circ\mu_C$... actually this requires
$h\colon TC\to TD$.  Instead we use:
\begin{align*}
    \mu_D\circ Th\circ\mu_C
    &= \mu_D\circ\mu_{TD}\circ T^2h \quad\text{(naturality of $\mu$ at $h$: applied as $\mu_D\circ Th = \ldots$)}
\end{align*}
Let us argue differently.  By naturality of~$\mu$ applied to~$Th\colon TC\to T^2D$:
\[
    \mu_D\circ T\mu_D = \mu_D\circ\mu_{TD}
    \quad\text{(monad associativity)}.
\]
And by naturality of~$\mu$ at the morphism $h\colon C\to TD$:
$\mu_{TD}\circ T^2 h = Th\circ\mu_C$... No, $h\colon C\to TD$,
so $Th\colon TC\to T^2D$ and $T^2h\colon T^2C\to T^3D$.
Naturality of $\mu$ gives $\mu_D\circ Th = h^*\circ\mu_C$... this is
getting confused.

The cleanest argument: define $f^{\sharp} = \mu_B\circ Tf$ for
$f\colon A\to TB$, so $g\circ_T f = g^{\sharp}\circ f$.  Then
associativity $(h\circ_T g)\circ_T f = h\circ_T(g\circ_T f)$
is equivalent to $(h^{\sharp}\circ g)^{\sharp} = h^{\sharp}\circ g^{\sharp}$,
i.e.\ $\mu\circ T(h^{\sharp}\circ g) = h^{\sharp}\circ\mu\circ Tg$,
which is $\mu\circ Th^{\sharp}\circ Tg = h^{\sharp}\circ\mu\circ Tg$.
This reduces to $\mu\circ Th^{\sharp} = h^{\sharp}\circ\mu$,
i.e.\ $\mu\circ T(\mu\circ Th) = \mu\circ Th\circ\mu$, which follows
from monad associativity and naturality of~$\mu$.
\end{proof}

\begin{proposition}\label{prop:kleisli-adjunction}
\index{Kleisli category!adjunction}
There is an adjunction $F_T\dashv G_T$ with
$F_T\colon\mathcal{C}\to\mathcal{C}_T$ and
$G_T\colon\mathcal{C}_T\to\mathcal{C}$ such that $G_TF_T = T$
and the monad arising from $F_T\dashv G_T$ is~$(T,\eta,\mu)$.
Explicitly:
\begin{itemize}
\item $F_T(A) = A$ on objects, $F_T(f) = \eta_B\circ f$ on morphisms
    $f\colon A\to B$,
\item $G_T(A) = TA$ on objects, $G_T(g) = \mu_B\circ Tg$ for
    $g\colon A\to TB$ (a Kleisli morphism $A\to B$).
\end{itemize}
\end{proposition}

\begin{proof}
The adjunction is verified directly:
$\Hom_{\mathcal{C}_T}(F_TA,B) = \Hom_{\mathcal{C}}(A,TB) = \Hom_{\mathcal{C}}(A,G_TB)$.
The unit is~$\eta$, the counit at~$B$ is $\id_{TB}\colon TB\to TB$
viewed as a Kleisli morphism $F_TG_TB = F_T(TB)\to B$.
One checks that $G_TF_T = T$ and the induced monad is $(T,\eta,\mu)$.
\end{proof}

\begin{remark}\label{rem:kleisli-em-universal}
\index{Kleisli category!as initial resolution}
\index{Eilenberg--Moore category!as terminal resolution}
The Kleisli category is the \textbf{initial} object and the
Eilenberg--Moore category is the \textbf{terminal} object in the
category of adjunctions giving rise to the monad~$T$.
More precisely, if $F\dashv G\colon\mathcal{D}\to\mathcal{C}$ is
any adjunction with associated monad~$T$, then there exist unique
(up to isomorphism) comparison functors
\[
    \mathcal{C}_T \longrightarrow \mathcal{D}
    \longrightarrow \mathcal{C}^T
\]
compatible with the adjunctions.  The comparison functor
$K\colon\mathcal{D}\to\mathcal{C}^T$ sends $D$ to $(GD,G\varepsilon_D)$.
\end{remark}

\begin{example}[Kleisli for the power set monad]\label{ex:kleisli-powerset}
\index{Kleisli category!power set monad}
\index{relation!as Kleisli morphism}
The Kleisli category for the power set monad~$\mathcal{P}$ on $\Set$
has sets as objects and $\Hom(A,B) = \Hom(A,\mathcal{P}(B)) \cong
\{R\subseteq A\times B\}$ as morphisms, i.e.\ morphisms are
\textbf{relations}.  Kleisli composition corresponds to composition
of relations.
\end{example}

\begin{example}[Kleisli for the Maybe monad]\label{ex:kleisli-maybe}
\index{Kleisli category!maybe monad}
The Kleisli category for the Maybe monad $T(X)=X\sqcup\{*\}$ has
morphisms $A\to B\sqcup\{*\}$, i.e.\ partial functions.  This
captures the idea of computations that may fail.
\end{example}


%% ============================================================
\section{Beck's monadicity theorem}\label{sec:beck}
%% ============================================================

The central question is: given an adjunction $F\dashv G$, when is the
comparison functor $K\colon\mathcal{D}\to\mathcal{C}^T$ an equivalence?

\begin{definition}[Comparison functor]\label{def:comparison-functor}
\index{comparison functor}
\index{$K$}
Let $F\dashv G\colon\mathcal{D}\to\mathcal{C}$ be an adjunction with
associated monad $(T,\eta,\mu)$ on~$\mathcal{C}$ (where $T = GF$).
The \textbf{comparison functor}
$K\colon\mathcal{D}\to\mathcal{C}^T$
is defined by:
\begin{itemize}
\item On objects: $K(D) = (GD,\, G\varepsilon_D)$, where
    $\varepsilon_D\colon FGD\to D$ is the counit.
    (This is a $T$-algebra: the unit and associativity axioms follow
    from the triangle identities and naturality of~$\varepsilon$.)
\item On morphisms: $K(f) = Gf$ for $f\colon D\to D'$.
    (This is a $T$-algebra morphism by naturality of~$\varepsilon$.)
\end{itemize}
\end{definition}

\begin{definition}[Monadic functor]\label{def:monadic}
\index{monadic functor}
\index{monadic adjunction}
An adjunction $F\dashv G$ (equivalently, the right adjoint~$G$) is
\textbf{monadic} if the comparison functor $K$ is an equivalence of
categories.  It is \textbf{strictly monadic} if $K$ is an isomorphism
of categories.
\end{definition}

\begin{definition}[$G$-split coequaliser]\label{def:g-split-coeq}
\index{coequaliser!$G$-split}
\index{split coequaliser}
\index{contractible coequaliser}
Let $G\colon\mathcal{D}\to\mathcal{C}$ be a functor.
A parallel pair $f,g\colon D\rightrightarrows D'$ in~$\mathcal{D}$ is
called a \textbf{$G$-split pair} if the image
$Gf,Gg\colon GD\rightrightarrows GD'$ admits a \textbf{split coequaliser}
in~$\mathcal{C}$: a diagram
\[
\begin{tikzcd}
GD \arrow[r, shift left, "Gf"] \arrow[r, shift right, "Gg"'] &
GD' \arrow[r, "e"] \arrow[l, bend left=30, "s"'] &
C \arrow[l, bend left=30, "t"']
\end{tikzcd}
\]
with $e\circ Gf = e\circ Gg$, $e\circ s = \id_C$, $Gf\circ t = s\circ e$,
and $Gg\circ t = \id_{GD'}$.

A \textbf{$G$-split coequaliser} of $f,g$ is a coequaliser of $f,g$
in~$\mathcal{D}$ whose image under~$G$ is this split coequaliser.
\end{definition}

\begin{remark}\label{rem:split-coeq-absolute}
\index{split coequaliser!absolute}
Split coequalisers are \textbf{absolute}: they are preserved by any
functor whatsoever.  Indeed, if $(e,s,t)$ is a split coequaliser as
above, then applying any functor~$H$ preserves all the relevant
equations, so $He$ is the coequaliser of $HGf,HGg$.  This crucial
property is what makes Beck's theorem work.
\end{remark}

\begin{lemma}\label{lem:split-coeq}
\index{split coequaliser!is a coequaliser}
A split coequaliser is a coequaliser.
\end{lemma}

\begin{proof}
Let $Gf,Gg\colon A\rightrightarrows B$ with split coequaliser $(e,s,t)$
as in Definition~\ref{def:g-split-coeq} (with $A=GD$, $B=GD'$).
First, $e$ coequalises: $e\circ Gf = e\circ Gg$ by assumption.
For universality, let $h\colon B\to Z$ with $h\circ Gf = h\circ Gg$.
Define $u = h\circ t\colon C\to Z$.  Then
$u\circ e = h\circ t\circ e = h\circ Gf\circ s = h\circ Gg\circ s$...
wait, we need: $u\circ e = h\circ t\circ e$.  Using $Gf\circ t = s\circ e$:
$h\circ s\circ e = h\circ Gf\circ t$... let us recompute.

We have $e\circ s = \id_C$, so $u\circ e = h\circ t\circ e$.
Now $t\circ e = Gf\circ t$... no, $Gf\circ t = s\circ e$.
Actually: from $Gg\circ t = \id_B$, we get
\[
    u\circ e = h\circ t\circ e = h\circ\id_B = h?
\]
No. Let us be careful.  We need:
From $Gg\circ t = \id_B$: this gives $t\colon C\to A$ is a section of
$Gg$... no, $t\colon C\to B$ and $Gg\colon A\to B$, so this does
not type-check.  Let me reconsider.

In the standard formulation, a split coequaliser of $f_1,f_2\colon A\rightrightarrows B$
is a diagram $A\rightrightarrows B\xrightarrow{e} C$ with
$s\colon C\to B$ and $t\colon B\to A$ such that
$e\circ f_1 = e\circ f_2$, $e\circ s = \id_C$,
$f_1\circ t = s\circ e$, $f_2\circ t = \id_B$.

Then for universality: let $h\colon B\to Z$ with $h\circ f_1 = h\circ f_2$.
Define $u = h\circ s\colon C\to Z$.  Then:
$u\circ e = h\circ s\circ e = h\circ f_1\circ t = h\circ f_2\circ t
= h\circ\id_B^{} = h$... wait: $f_2\circ t = \id_B$, so
$h\circ f_2\circ t = h$.  And $h\circ f_1\circ t = h\circ f_2\circ t = h$
since $h\circ f_1 = h\circ f_2$.  So $u\circ e = h\circ s\circ e
= h\circ f_1\circ t = h$.

Uniqueness: if $u'\circ e = h$, then $u' = u'\circ e\circ s = h\circ s = u$.
\end{proof}

Now we can state Beck's theorem.

\begin{theorem}[Beck's Monadicity Theorem]\label{thm:beck}
\index{Beck's monadicity theorem}
\index{monadicity theorem}
\index{monadic functor!characterisation}
Let $F\dashv G$ with $G\colon\mathcal{D}\to\mathcal{C}$.
The following are equivalent:
\begin{enumerate}[label=(\roman*)]
\item $G$ is monadic (the comparison functor $K$ is an equivalence).
\item $G$ reflects isomorphisms and $\mathcal{D}$ has coequalisers of
    all $G$-split pairs, and $G$ preserves them.
\end{enumerate}
\end{theorem}

\begin{proof}
The proof is substantial.  We break it into parts.

\medskip\noindent
\textbf{(i) $\Rightarrow$ (ii):}

Assume $K\colon\mathcal{D}\to\mathcal{C}^T$ is an equivalence.
Since $U^T\circ K = G$ and $K$ is an equivalence, it suffices to
show that $U^T\colon\mathcal{C}^T\to\mathcal{C}$ reflects isomorphisms
and that $\mathcal{C}^T$ has coequalisers of $U^T$-split pairs.

\emph{$U^T$ reflects isomorphisms:}
Let $f\colon(A,a)\to(B,b)$ be a $T$-algebra morphism such that
$f\colon A\to B$ is an isomorphism in~$\mathcal{C}$.
We claim $f^{-1}\colon(B,b)\to(A,a)$ is a $T$-algebra morphism.
We need $f^{-1}\circ b = a\circ Tf^{-1}$.  Since $f\circ a = b\circ Tf$
(as $f$ is a $T$-algebra morphism), we get
$a = f^{-1}\circ b\circ Tf$, hence
$a\circ Tf^{-1} = f^{-1}\circ b\circ Tf\circ Tf^{-1}
= f^{-1}\circ b\circ T(f\circ f^{-1}) = f^{-1}\circ b$.

\emph{$\mathcal{C}^T$ has coequalisers of $U^T$-split pairs:}
Let $f,g\colon(A,a)\rightrightarrows(B,b)$ be $T$-algebra morphisms
such that $f,g\colon A\rightrightarrows B$ admits a split coequaliser
$(e\colon B\to C,\, s\colon C\to B,\, t\colon B\to A)$ in~$\mathcal{C}$.
We must endow~$C$ with a $T$-algebra structure.

Define $c\colon TC\to C$ by
\[
    c \;=\; e\circ b\circ Ts \;\colon\; TC\to C.
\]
\emph{$c$ is well-defined as a structure map:}
We verify the algebra axioms.  Unit: $c\circ\eta_C = e\circ b\circ Ts\circ\eta_C
= e\circ b\circ\eta_B\circ s = e\circ s = \id_C$
(using naturality of~$\eta$ and the algebra unit for~$(B,b)$).
Associativity: $c\circ\mu_C = e\circ b\circ Ts\circ\mu_C
= e\circ b\circ T(b\circ Ts) = e\circ b\circ Tb\circ T^2s$...
We use the algebra associativity for $(B,b)$: $b\circ Tb = b\circ\mu_B$.
So $c\circ\mu_C = e\circ b\circ\mu_B\circ T^2s = e\circ b\circ Ts\circ T(e\circ b\circ Ts)$...
let us use $\mu_C = T(s)\circ... $ no. Actually:
\begin{align*}
    c\circ Tc &= (e\circ b\circ Ts)\circ T(e\circ b\circ Ts) \\
    &= e\circ b\circ Ts\circ Te\circ Tb\circ T^2s \\
    &= e\circ b\circ T(s\circ e)\circ Tb\circ T^2s \\
    &= e\circ b\circ T(f\circ t)\circ Tb\circ T^2s
    \quad\text{(since $s\circ e = f\circ t$... no, $f_1\circ t = s\circ e$)}
\end{align*}
We have $f\circ t = s\circ e$ from the split coequaliser.
So $T(s\circ e) = Tf\circ Tt$.  Then:
\begin{align*}
    c\circ Tc &= e\circ b\circ Tf\circ Tt\circ Tb\circ T^2s \\
    &= e\circ f\circ a\circ Tt\circ Tb\circ T^2s
    \quad\text{($b\circ Tf = f\circ a$ since $f$ is a $T$-morphism)}\\
    &= e\circ g\circ a\circ Tt\circ Tb\circ T^2s
    \quad\text{($e\circ f = e\circ g$)}\\
    &= e\circ b\circ Tg\circ Tt\circ Tb\circ T^2s
    \quad\text{($b\circ Tg = g\circ a$ since $g$ is a $T$-morphism)}\\
    &= e\circ b\circ T(g\circ t)\circ Tb\circ T^2s \\
    &= e\circ b\circ T(\id_B)\circ Tb\circ T^2s
    \quad\text{($g\circ t = \id_B$)} \\
    &= e\circ b\circ Tb\circ T^2s \\
    &= e\circ b\circ\mu_B\circ T^2s
    \quad\text{(associativity for $(B,b)$)} \\
    &= e\circ b\circ Ts\circ\mu_C
    \quad\text{(naturality of $\mu$ at $s$: $\mu_C = Ts... $ no)}
\end{align*}
Naturality of~$\mu$ at $s\colon C\to B$ gives
$\mu_B\circ T^2s = Ts\circ\mu_C$.  Hence:
\[
    c\circ Tc = e\circ b\circ Ts\circ\mu_C = c\circ\mu_C.
\]
So $(C,c)$ is indeed a $T$-algebra.

\emph{$e$ is a $T$-algebra morphism:}
We check $e\circ b = c\circ Te$.
$c\circ Te = e\circ b\circ Ts\circ Te = e\circ b\circ T(s\circ e)
= e\circ b\circ Tf\circ Tt = e\circ f\circ a\circ Tt$.
But also $e\circ b = e\circ b\circ T(g\circ t) = e\circ g\circ a\circ Tt
= e\circ f\circ a\circ Tt$... this is not right either.
We need $c\circ Te = e\circ b$ directly.
$c\circ Te = e\circ b\circ T(s\circ e) = e\circ b\circ T(f\circ t)
= e\circ b\circ Tf\circ Tt = e\circ f\circ a\circ Tt$.
And $e\circ b = e\circ b\circ\id_{TB} = e\circ b\circ T(g\circ t)
= e\circ g\circ a\circ Tt = e\circ f\circ a\circ Tt$.
Here we used $g\circ t = \id_B$ and $e\circ g = e\circ f$.
So indeed $c\circ Te = e\circ b$.

\emph{$(C,c)$ is the coequaliser:}
Suppose $h\colon(B,b)\to(Z,z)$ is a $T$-algebra morphism with
$h\circ f = h\circ g$.  Since the underlying diagram is a split
coequaliser, there is a unique $u\colon C\to Z$ in~$\mathcal{C}$ with
$u\circ e = h$.  We must show $u$ is a $T$-algebra morphism, i.e.\
$u\circ c = z\circ Tu$.
\begin{align*}
    u\circ c &= u\circ e\circ b\circ Ts = h\circ b\circ Ts
    = z\circ Th\circ Ts = z\circ T(h\circ s)
    = z\circ T(u\circ e\circ s) = z\circ Tu.
\end{align*}
(Using $h\circ b = z\circ Th$ since $h$ is a $T$-morphism, and
$e\circ s = \id_C$.)

\medskip\noindent
\textbf{(ii) $\Rightarrow$ (i):}

Assume $G$ reflects isomorphisms and $\mathcal{D}$ has coequalisers
of $G$-split pairs, preserved by~$G$.  We show $K$ is an equivalence
by constructing a quasi-inverse.

\emph{$K$ is essentially surjective:}
Let $(A,a)\in\mathcal{C}^T$.  Consider the pair
$FTA \rightrightarrows FA$ in~$\mathcal{D}$ given by
\[
    Fa \;\text{ and }\; F\mu_A \quad (\text{no: }
    \varepsilon_{FA}\colon FGFA = FTA\to FA).
\]
More precisely, consider the canonical pair:
$F\mu_A, \varepsilon_{FTA}\colon FT^2A = FGFTA\rightrightarrows FTA$...

Let us use the standard construction.  For $(A,a)\in\mathcal{C}^T$,
consider the parallel pair in~$\mathcal{D}$:
\[
    FTGA \;=\; FTA \xrightrightarrows[\varepsilon_{FA}]{Fa} FA.
\]
Wait---$a\colon TA\to A$ is a morphism in~$\mathcal{C}$, so $Fa\colon FTA\to FA$.
And $\varepsilon_{FA}\colon FGFA = FTA\to FA$.  So we have
$Fa,\varepsilon_{FA}\colon FTA\rightrightarrows FA$.

Applying~$G$:
$GFa,G\varepsilon_{FA}\colon GFTA = T^2A\rightrightarrows TA = GFA$.
That is, $Ta,\mu_A\colon T^2A\rightrightarrows TA$.
This pair has a split coequaliser in~$\mathcal{C}$ with
$e = a\colon TA\to A$, $s = \eta_A\colon A\to TA$,
$t = T\eta_A... $ let us verify:
\begin{itemize}
\item $e\circ Ta = a\circ Ta$ and $e\circ\mu_A = a\circ\mu_A$.
    The algebra associativity says $a\circ Ta = a\circ\mu_A$. \checkmark
\item $e\circ s = a\circ\eta_A = \id_A$. \checkmark\ (algebra unit)
\item $Ta\circ t = s\circ e$: We need $t\colon TA\to T^2A$ with
    $Ta\circ t = \eta_A\circ a$ and $\mu_A\circ t = \id_{TA}$.
    Take $t = T\eta_A$.  Then $\mu_A\circ T\eta_A = \id_{TA}$
    (monad right unit). \checkmark
\item $Ta\circ T\eta_A = T(a\circ\eta_A) = T(\id_A) = \id_{TA}$.
    But we need $Ta\circ t = s\circ e = \eta_A\circ a$.
    We have $Ta\circ T\eta_A = \id_{TA}\neq\eta_A\circ a$ in general.
\end{itemize}

So let us take $t = \eta_{TA}\colon TA\to T^2A$ instead.
\begin{itemize}
\item $\mu_A\circ\eta_{TA} = \id_{TA}$. \checkmark\ (monad left unit)
\item $Ta\circ\eta_{TA} = \eta_A\circ a$: this is naturality of~$\eta$
    at~$a$. \checkmark
\end{itemize}
So the split coequaliser data is: $e = a$, $s = \eta_A$, $t = \eta_{TA}$,
with the coequalised pair being $(Ta,\mu_A)$.

So the pair $(Fa,\varepsilon_{FA})$ is a $G$-split pair.
By hypothesis, $\mathcal{D}$ has a coequaliser
$q\colon FA\to D$ of $(Fa,\varepsilon_{FA})$, and $G$ preserves it.
Since $G$ preserves this coequaliser and the split coequaliser of
$(Ta,\mu_A)$ is $a\colon TA\to A$, we get $Gq = a$ (up to the
canonical isomorphism).  More precisely, $GD\cong A$ via the
isomorphism coming from the fact that $Gq\colon GFA = TA\to GD$
is the coequaliser of $(Ta,\mu_A)$ in~$\mathcal{C}$, which is~$a$
with target~$A$.  So $GD\cong A$.

Now $K(D) = (GD, G\varepsilon_D)$.
Under the identification $GD\cong A$, we get a $T$-algebra structure
on~$A$ that agrees with~$a$.  Hence $K(D)\cong(A,a)$, and $K$ is
essentially surjective.

\emph{$K$ is fully faithful:}
For $D,D'\in\mathcal{D}$, the map
$K_{D,D'}\colon\Hom_{\mathcal{D}}(D,D')\to
\Hom_{\mathcal{C}^T}(K(D),K(D'))$
sends $f\mapsto Gf$.  Since morphisms of $T$-algebras are just
morphisms in~$\mathcal{C}$ satisfying a compatibility condition,
$\Hom_{\mathcal{C}^T}(K(D),K(D')) \subseteq \Hom_{\mathcal{C}}(GD,GD')$.

\emph{Faithful:}
If $Gf = Gf'$, we want $f = f'$.  Consider the coequalisers
$q\colon FA\to D$ from the essential surjectivity argument.
Since $Gf = Gf'$ and $G$ reflects the relevant coequalisers, one
shows $f = f'$.  Alternatively: for any $D\in\mathcal{D}$,
the counit $\varepsilon_D\colon FGD\to D$ is a (regular) epimorphism
(it is the coequaliser of a $G$-split pair).
If $Gf = Gf'$, then $FGf = FGf'$, so
$f\circ\varepsilon_D = \varepsilon_{D'}\circ FGf
= \varepsilon_{D'}\circ FGf' = f'\circ\varepsilon_D$.
Since $\varepsilon_D$ is an epimorphism, $f = f'$.

\emph{Full:}
Let $\varphi\colon GD\to GD'$ be a $T$-algebra morphism
$K(D)\to K(D')$.  Consider the diagram in~$\mathcal{D}$:
\[
\begin{tikzcd}
FGFGD \arrow[r, shift left, "F(G\varepsilon_D)"]
       \arrow[r, shift right, "\varepsilon_{FGD}"'] &
FGD \arrow[r, "\varepsilon_D"] &
D
\end{tikzcd}
\]
and the corresponding one for~$D'$.  The map $F\varphi\colon FGD\to FGD'$
satisfies the required compatibility (since $\varphi$ is a $T$-algebra
morphism).  Since $\varepsilon_D$ is the coequaliser and
$\varepsilon_{D'}\circ F\varphi$ coequalises the pair for~$D$,
there exists a unique $f\colon D\to D'$ with
$f\circ\varepsilon_D = \varepsilon_{D'}\circ F\varphi$.
Applying~$G$: $Gf\circ G\varepsilon_D = G\varepsilon_{D'}\circ GF\varphi
= G\varepsilon_{D'}\circ T\varphi$.  But for a $T$-algebra morphism,
$\varphi\circ G\varepsilon_D = G\varepsilon_{D'}\circ T\varphi$, so
$Gf\circ G\varepsilon_D = \varphi\circ G\varepsilon_D$.
Since $G\varepsilon_D = a$ is an epimorphism (it is a split epimorphism),
$Gf = \varphi$.

Hence $K$ is fully faithful and essentially surjective, so it is an
equivalence.
\end{proof}

\begin{remark}\label{rem:beck-crude}
\index{crude monadicity theorem}
There is a \textbf{crude monadicity theorem} (also called the
\textbf{tripleability theorem} in older terminology): $G$ is monadic
if and only if $G$ has a left adjoint, $G$ reflects isomorphisms,
and $\mathcal{D}$ has and $G$ preserves coequalisers of
\emph{reflexive} $G$-split pairs (a reflexive pair is one admitting
a common section).  This is sometimes easier to verify in practice.
\end{remark}

\begin{corollary}\label{cor:monadic-reflects}
\index{monadic functor!reflects isomorphisms}
A monadic functor reflects isomorphisms and creates coequalisers of
$G$-split pairs.
\end{corollary}

\begin{proof}
This follows immediately from the (i) $\Rightarrow$ (ii) direction
of Beck's theorem.
\end{proof}


%% ============================================================
\section{Applications of monadicity}\label{sec:monadicity-apps}
%% ============================================================

\begin{example}[Algebraic categories are monadic over $\Set$]%
\label{ex:monadic-algebraic}
\index{monadic functor!algebraic categories}
\index{algebraic category}
The forgetful functors from $\Grp$, $\Ab$, $\mathsf{Ring}$,
$R\text{-}\mathsf{Mod}$, and more generally any variety of
algebras in the sense of universal algebra, to~$\Set$ are monadic.
We verify the conditions of Beck's theorem:
\begin{enumerate}[label=(\alph*)]
\item The forgetful functor $U$ has a left adjoint (the free functor).
\item $U$ reflects isomorphisms: a bijective homomorphism is an
    isomorphism.
\item $\Set$ has all coequalisers, and the free-forgetful adjunction
    creates them via the algebra structure.
\end{enumerate}
\end{example}

\begin{example}[Compact Hausdorff spaces]\label{ex:monadic-comphaus}
\index{monadic functor!compact Hausdorff spaces}
\index{compact Hausdorff space}
The forgetful functor $U\colon\mathsf{CompHaus}\to\Set$ is monadic.
The associated monad is the ultrafilter monad~$\beta$, whose
algebras are compact Hausdorff spaces.  This is a deep theorem of
Manes (1969).
\index{ultrafilter monad}
\index{Manes' theorem}
\end{example}

\begin{example}[Non-monadic functors]\label{ex:non-monadic}
\index{monadic functor!non-examples}
\begin{enumerate}[label=(\roman*)]
\item The forgetful functor $U\colon\Top\to\Set$ is \emph{not} monadic.
    It has a left adjoint (discrete topology) but does not reflect
    isomorphisms: a continuous bijection need not be a homeomorphism.

\item The inclusion $\Ab\hookrightarrow\Grp$ is not monadic:
    it does not reflect isomorphisms in general, and the associated
    monad (abelianisation) does not recover~$\Ab$ as the category of
    algebras (the Eilenberg--Moore category is equivalent to~$\Ab$,
    but the inclusion is not the comparison functor from~$\Grp$).
    Actually, the functor $\Ab\to\Grp$ does reflect isomorphisms
    (it is fully faithful), so the issue is that coequalisers of
    split pairs in~$\Grp$ need not land in~$\Ab$.
\end{enumerate}
\end{example}

\begin{proposition}\label{prop:monadic-creates-limits}
\index{monadic functor!creates limits}
A monadic functor creates all limits that exist in the base category.
In particular, if $G\colon\mathcal{D}\to\mathcal{C}$ is monadic
and $\mathcal{C}$ is complete, then so is~$\mathcal{D}$.
\end{proposition}

\begin{proof}
Since $G$ is monadic, $\mathcal{D}\simeq\mathcal{C}^T$ for a
monad~$T$ on~$\mathcal{C}$.  It suffices to show $U^T$ creates limits.
Let $D\colon\mathcal{J}\to\mathcal{C}^T$ be a diagram with
$D(j) = (A_j,a_j)$, and suppose $L = \varprojlim_j A_j$ exists
in~$\mathcal{C}$ with projections $\pi_j\colon L\to A_j$.
Then $TL = T(\varprojlim A_j)$ and we need a $T$-algebra structure
on~$L$.  Since $T$ (being a right adjoint composite $G\circ F$... no,
$T$ need not preserve limits in general).

For each~$j$, the maps $a_j\circ T\pi_j\colon TL\to A_j$ form a cone
over the diagram $\{A_j\}$ (using compatibility of the $a_j$ with
the algebra morphisms).  By the universal property of~$L$, there is a
unique $\ell\colon TL\to L$ with $\pi_j\circ\ell = a_j\circ T\pi_j$
for all~$j$.  One verifies that $(L,\ell)$ is a $T$-algebra and is
the limit of~$D$ in~$\mathcal{C}^T$.  The algebra axioms for~$(L,\ell)$
follow from those for each~$(A_j,a_j)$ and the fact that the $\pi_j$
jointly detect equality (since $L$ is their limit).
\end{proof}


%% ============================================================
\section{Exercises}\label{sec:monad-exercises}
%% ============================================================

\begin{exercise}\label{exer:monad-triangle}
\index{monad!from adjunction}
Complete the details of Proposition~\ref{prop:adj-monad}:
verify each monad axiom carefully using the triangle identities.
\end{exercise}

\begin{exercise}\label{exer:monad-list}
\index{monad!list monad}
\index{list monad}
Let $T\colon\Set\to\Set$ be the ``list'' functor with
$T(X) = \coprod_{n\geq 0}X^n$.  Define $\eta$ and $\mu$ explicitly
and verify the monad axioms.  Show that $T$-algebras are monoids.
\end{exercise}

\begin{exercise}\label{exer:em-adjoint}
\index{Eilenberg--Moore category!forgetful-free adjunction}
Verify all details of Proposition~\ref{prop:em-adjunction}: show that
$F^T\dashv U^T$ is an adjunction and that the associated monad is
$(T,\eta,\mu)$.
\end{exercise}

\begin{exercise}\label{exer:kleisli-composition}
\index{Kleisli category!composition}
For the free group monad $T$ on $\Set$, describe Kleisli composition
explicitly.  Show that a Kleisli morphism $A\to B$ is a function
$A\to F(B)$ (where $F(B)$ is the free group on~$B$) and that
composition corresponds to substitution followed by reduction.
\end{exercise}

\begin{exercise}\label{exer:idempotent-monad}
\index{monad!idempotent}
\index{idempotent monad}
A monad $(T,\eta,\mu)$ is \textbf{idempotent} if
$\mu\colon T^2\Rightarrow T$ is a natural isomorphism.
\begin{enumerate}[label=(\alph*)]
\item Show that the following are equivalent: (i) $T$ is idempotent;
    (ii) $T\eta = \eta T$; (iii) every $T$-algebra is free.
\item Show that the monad arising from a reflective subcategory
    (i.e.\ a full subcategory whose inclusion has a left adjoint)
    is idempotent.
\item Deduce that the category of algebras for an idempotent monad
    is equivalent to the full subcategory of~$\mathcal{C}$ consisting
    of objects of the form~$TA$.
\end{enumerate}
\end{exercise}

\begin{exercise}\label{exer:monad-distributive-law}
\index{monad!distributive law}
\index{distributive law}
Let $(S,\eta^S,\mu^S)$ and $(T,\eta^T,\mu^T)$ be monads on~$\mathcal{C}$.
A \textbf{distributive law} of $S$ over~$T$ is a natural transformation
$\lambda\colon TS\Rightarrow ST$ satisfying four axioms (compatibility
with units and multiplications of both monads).
\begin{enumerate}[label=(\alph*)]
\item Write down the four axiom diagrams.
\item Show that given a distributive law, the composite $ST$ carries a
    monad structure with unit $\eta^S\eta^T$ and multiplication
    $\mu^S\cdot S\lambda T\cdot\mu^T$.
\item Show that the ring axiom ``multiplication distributes over
    addition'' is a distributive law for the free monoid monad over
    the free abelian group monad.
\end{enumerate}
\end{exercise}

\begin{exercise}\label{exer:beck-groups}
\index{Beck's monadicity theorem!applied to groups}
Use Beck's monadicity theorem to give an alternative proof that the
forgetful functor $U\colon\Grp\to\Set$ is monadic.
\emph{Hint:} Verify that $U$ reflects isomorphisms and that $\Grp$
has coequalisers of $U$-split pairs.
\end{exercise}

\begin{exercise}\label{exer:monadic-top}
\index{monadic functor!$\Top$ is not monadic over $\Set$}
Show that the forgetful functor $U\colon\Top\to\Set$ is not monadic by
exhibiting a continuous bijection that is not a homeomorphism.
What is the monad $T = UF$ (where $F$ is the discrete topology functor)?
\end{exercise}

\begin{exercise}\label{exer:comonad-coalgebra}
\index{comonad!coalgebra}
\index{coalgebra!over a comonad}
Dualise the theory of this chapter: define a \textbf{coalgebra} for a
comonad $(S,\varepsilon,\delta)$ on~$\mathcal{C}$ and formulate the
dual of Beck's theorem for comonadic functors.
\end{exercise}

\begin{exercise}\label{exer:monad-maps}
\index{monad!morphism}
A \textbf{morphism of monads} $\varphi\colon(S,\eta^S,\mu^S)\to
(T,\eta^T,\mu^T)$ on the same category~$\mathcal{C}$ is a natural
transformation $\varphi\colon S\Rightarrow T$ compatible with units and
multiplications:
$\varphi\circ\eta^S = \eta^T$ and
$\varphi\circ\mu^S = \mu^T\circ\varphi\varphi$
(where $\varphi\varphi = \varphi T\circ S\varphi = T\varphi\circ\varphi S$).
\begin{enumerate}[label=(\alph*)]
\item Show that a monad morphism $\varphi\colon S\to T$ induces a functor
    $\varphi^*\colon\mathcal{C}^T\to\mathcal{C}^S$ (restriction of scalars).
\item Show that $\varphi^*$ has a left adjoint (extension of scalars).
\end{enumerate}
\end{exercise}

\begin{exercise}\label{exer:monad-coequaliser-detail}
\index{Beck's monadicity theorem!coequaliser construction}
In the proof of Beck's theorem (direction (i) $\Rightarrow$ (ii)),
verify directly that the candidate $T$-algebra structure
$c = e\circ b\circ Ts$ on~$C$ satisfies the associativity axiom, filling
in any steps that were left implicit.
\end{exercise}

\index{monad|)}
